\chapter{Einleitung}
\label{ch:einleitung}

Die Soziale Arbeit steht vor der Herausforderung, ihre Methoden kontinuierlich an gesellschaftliche Veränderungen anzupassen \parencite[vgl.][12\psq]{mueller2020}. Dabei spielt die Professionalisierung eine zentrale Rolle. Wie \textcite[42]{bonhoeffer1994} betont: \enquote{Die Verantwortung des Einzelnen kann nicht an Institutionen delegiert werden.}

Im Kontext der \gls{llm}-gestützten Analyse ergeben sich neue Möglichkeiten für die empirische Sozialforschung. Die \gls{api}-basierte Datenerhebung ermöglicht es, große Textmengen systematisch auszuwerten \parencite[vgl.][]{schmidt2019}.

\section{Problemstellung}

Die vorliegende Arbeit befasst sich mit der Frage, wie digitale Methoden in der Sozialen Arbeit eingesetzt werden können. Insbesondere wird untersucht, welche Rolle \gls{nlp} bei der Auswertung qualitativer Daten spielen kann.

\section{Zielsetzung und Aufbau der Arbeit}

Ziel dieser Bachelorarbeit ist es, einen Überblick über aktuelle Ansätze zu geben und diese kritisch zu bewerten. Die Arbeit gliedert sich in folgende Abschnitte:

\begin{enumerate}
    \item Theoretischer Hintergrund
    \item Methodisches Vorgehen
    \item Ergebnisse und Diskussion
    \item Fazit und Ausblick
\end{enumerate}
